\documentclass[output=paper,colorlinks,citecolor=brown]{langscibook}
\ChapterDOI{10.5281/zenodo.20541852}
\author{Dominika Skrzypek\orcid{}\affiliation{Adam Mickiewicz University Poznań}}

\title[The future that did not happen]{The future that did not happen: On the history of the verb \textit{varda} in Swedish}

\abstract{This paper considers the constructions with the verb \textit{varda} ‘to become’ in the history of the Swedish language, in particular such uses of the verb (with infinitives, present participles or in the present tense) which could constitute future reference. We follow the distribution and usage of each construction across a time span between 1300 and 1750 and complete with a qualitative analysis of the extant examples. We also consider the reasons for the demise of the futurate construction with the infinitive, in particular its low frequencies, limited distribution and the role played by the rise of a passive construction.}

%move the following commands to the "local..." files of the master project when integrating this chapter

% \usepackage{multirow}
% \usepackage{hanging}


%let\eachwordone=\itshape

\IfFileExists{../localcommands.tex}{
   \addbibresource{../localbibliography.bib}
   \usepackage{tabularx,multicol}
\usepackage{url}
\urlstyle{same}

\usepackage[english]{babel}
 
\usepackage{langsci-optional}
\usepackage{langsci-lgr}
\usepackage{langsci-gb4e}

\usepackage{langsci-branding}

%\usepackage{tabularx}

% elena
\usepackage{graphicx} %package to manage images
\usepackage{paralist}

% sarah & lena
\usepackage[justification=centering]{caption}
\usepackage{rotating}

% robin
\usepackage{lscape}
\usepackage{tikz}
\usetikzlibrary{tikzmark}

% dominika
\usepackage{multirow}
\usepackage{hanging}

% olaf & dylan
\usepackage{subcaption}

% svetlana
\usepackage{verbatim}
\usepackage{enumitem}

% wiemer
%\usepackage[markup=underlined]{changes}
%\usepackage{hyphenat}

% wir
\usepackage{longtable}

% for crossref in the introduction jiabin
\usepackage{xr}
   % \newcommand*{\orcid}{}
%

\makeatletter
\let\thetitle\@title
\let\theauthor\@author
\makeatother

\newcommand{\togglepaper}[1][0]{
   \bibliography{../localbibliography}
   \papernote{\scriptsize\normalfont
     \theauthor.
     \titleTemp.
     To appear in:
     E. Di Tor \& Herr Rausgeberin (ed.).
     Booktitle in localcommands.tex.
     Berlin: Language Science Press. [preliminary page numbering]
   }
   \pagenumbering{roman}
   \setcounter{chapter}{#1}
   \addtocounter{chapter}{-1}
}
%
\newbool{bookcompile}
\booltrue{bookcompile}
\newcommand{\bookorchapter}[2]{\ifbool{bookcompile}{#1}{#2}}
%
% %\newcommand{\orcid}[1]{}
\graphicspath{ {./figures/} }
%
\let\eachwordone=\itshape


% Seongha
% \newcommand{\koreex}[1]{{\fontspec{Baekmuk Batang}{#1}}}

% Kutida
\newfontfamily\thaifont{Noto Sans Thai} 
\newcommand{\textthai}[1]{{\thaifont #1}} % für Thai


\newfontfamily\chinesefont{Noto Sans CJK SC}
\newcommand{\textchinese}[1]{{\chinesefont #1}} % Für Chinesisch

% wiemer
\def\hyph{-\penalty0\hskip0pt\relax}
% \definechangesauthor[color=NavyBlue]{SH}



\newcolumntype{K}{>{\raggedright\arraybackslash}p{0.6\textwidth}}  % Right column with 0.7 width


% \newfontfamily\krn[Scale=MatchLowercase,BoldFont=SourceHanSerif-Bold.otf]{SourceHanSerif-Regular.otf}
% \AdditionalFontImprint{Source Han Serif KO}
% \XeTeXlinebreaklocale 'ko'
% \renewcommand{\krn}{}



% Rhee
\newfontfamily\koreanfont{Noto Sans CJK KR} 
\newcommand{\textkorean}[1]{{\koreanfont #1}} % für Koreanisch

\renewcommand{\REF}[2][]{(\ref{#2}#1)}
   %% hyphenation points for line breaks
%% Normally, automatic hyphenation in LaTeX is very good
%% If a word is mis-hyphenated, add it to this file
%%
%% add information to TeX file before \begin{document} with:
%% %% hyphenation points for line breaks
%% Normally, automatic hyphenation in LaTeX is very good
%% If a word is mis-hyphenated, add it to this file
%%
%% add information to TeX file before \begin{document} with:
%% %% hyphenation points for line breaks
%% Normally, automatic hyphenation in LaTeX is very good
%% If a word is mis-hyphenated, add it to this file
%%
%% add information to TeX file before \begin{document} with:
%% \include{localhyphenation}

\hyphenation{
    con-tri-bute
    par-a-digm
    pe-ri-phra-ses
    achieve-ment
    accomp-lish-ment
    We-li-nchen-kang-ci-kok
    cross-ling-u-is-tic
    rec-ord-s
    trun-ca-ted
    de-ri-va-tion-al
    stem-de-ri-va-tion-al
    pre-sent
    in-flu-en-ced
    sys-tem
    ques-tions
    Bor-kovs-kij
    Da-ny-len-ko
    e-qui-va-lents
    re-cog-nized
    In-ter-fe-renz-schicht
    Aus-ei-nan-der-se-tz-ung
    Früh-neu-hoch-deutsch-kor-pus
    Ak-tions-art
    Smir-no-va
    Whi-le
    Ar-chive
    bi-ginnan
    Rei-chen-bach
    ger-ma-nis-ti-sche
    Re-fe-renz-kor-pus
    le-xi-ka-lisch-en
    Alt-hoch-deutsche
    Li-te-ra-tur
    dia-chron
    Dia-lek-to-lo-gie
}

\hyphenation{
    con-tri-bute
    par-a-digm
    pe-ri-phra-ses
    achieve-ment
    accomp-lish-ment
    We-li-nchen-kang-ci-kok
    cross-ling-u-is-tic
    rec-ord-s
    trun-ca-ted
    de-ri-va-tion-al
    stem-de-ri-va-tion-al
    pre-sent
    in-flu-en-ced
    sys-tem
    ques-tions
    Bor-kovs-kij
    Da-ny-len-ko
    e-qui-va-lents
    re-cog-nized
    In-ter-fe-renz-schicht
    Aus-ei-nan-der-se-tz-ung
    Früh-neu-hoch-deutsch-kor-pus
    Ak-tions-art
    Smir-no-va
    Whi-le
    Ar-chive
    bi-ginnan
    Rei-chen-bach
    ger-ma-nis-ti-sche
    Re-fe-renz-kor-pus
    le-xi-ka-lisch-en
    Alt-hoch-deutsche
    Li-te-ra-tur
    dia-chron
    Dia-lek-to-lo-gie
}

\hyphenation{
    con-tri-bute
    par-a-digm
    pe-ri-phra-ses
    achieve-ment
    accomp-lish-ment
    We-li-nchen-kang-ci-kok
    cross-ling-u-is-tic
    rec-ord-s
    trun-ca-ted
    de-ri-va-tion-al
    stem-de-ri-va-tion-al
    pre-sent
    in-flu-en-ced
    sys-tem
    ques-tions
    Bor-kovs-kij
    Da-ny-len-ko
    e-qui-va-lents
    re-cog-nized
    In-ter-fe-renz-schicht
    Aus-ei-nan-der-se-tz-ung
    Früh-neu-hoch-deutsch-kor-pus
    Ak-tions-art
    Smir-no-va
    Whi-le
    Ar-chive
    bi-ginnan
    Rei-chen-bach
    ger-ma-nis-ti-sche
    Re-fe-renz-kor-pus
    le-xi-ka-lisch-en
    Alt-hoch-deutsche
    Li-te-ra-tur
    dia-chron
    Dia-lek-to-lo-gie
}
   \boolfalse{bookcompile}
   \togglepaper[5]%%chapternumber
}{}

\begin{document}
\maketitle

\section{Introduction} \label{sec:skrzypek:1}
Typologically, future constructions develop from a relatively small set of lexical sources and follow restricted developmental paths. Cross-linguistically, the main lexical sources of future constructions are motion verbs such as ‘come’ and ‘go’ on the one hand and verbs of ability, obligation, and desire on the other hand (\cite{Heine1995}). The grammaticalization into future tense markers involves a convergence of these different sources and a development along similar paths, most notably including the notion of intention. As argued by \citet[254]{BybeePerkinsPagliuca1994} ‘all futures go through a stage of functioning to express the intention, first of the speaker, and later of the agent of the main verb.’ \citet{Heine1995} argues that even those future tense constructions which are derived from other sources than the ones mentioned above, e.g., the German \textit{werden} ‘become’, at one stage become markers of intention before proceeding along the grammaticalization cline (\cite[127]{Heine1995}). However, it has also been demonstrated that the development of de-venitive future constructions need not involve the notion of intentionality (\cite[322]{Dahl2000}), but rather the development proceeds through ingressive, inchoative or resultative meanings (\cite[378]{Traugott1978}). 

Apart from the lexically based futures (future constructions which grammaticalize from lexical sources), also aspectual forms can acquire future time reference, so that present tense forms can be used to refer to future events (so-called aspectual futures, see \citealt[275]{BybeePerkinsPagliuca1994}, or futurate present, see \citealt[157--180]{Hilpert2008}). However, in contrast to the lexically based futures, with the aspectual futures, future time reference is not part of their semantics, but is only achieved as a contextual effect. 

In Modern Swedish (MSw), there are a number of grammaticalized means of expressing futurity, which include: de-venitive \textit{komma att V}, lit. ‘come to V’, modal verbs, in particular \textit{ska} ‘shall’ and the use of the present tense, typically with temporal adverbials indicating its future reading. The constructions are exemplified in (\ref{ex:skrzypek:1}--\ref{ex:skrzypek:3}).

\ea \label{ex:skrzypek:1}
    \gll Jag kommer att resa imorgon.\\
    I come.\textsc{prs} to travel.\textsc{inf} tomorrow.\\
    \glt ‘I am going to travel tomorrow.’
\z

\ea \label{ex:skrzypek:2}
    \gll Jag ska resa imorgon.\\
    I shall.\textsc{prs} travel.\textsc{inf} tomorrow.\\
    \glt ‘I will travel tomorrow.’
\z

\ea \label{ex:skrzypek:3}
    \gll Jag reser imorgon.\\
    I travel.\textsc{prs} tomorrow.\\
    \glt ‘I travel tomorrow.’
\z

There is in MSw no construction corresponding to the German \textit{werden} construction, although a cognate of the verb, \textit{varda} ‘become’, is present in some Swedish dialects (in particular in the northern regions, \citet{Eklund2010}, \citet{Loenheim2013}). The verb is also well-attested in the oldest extant text from the Swedish territory, the runic inscriptions, and in the oldest Swedish prose written with Latin script, from early 1200s (\cite{Skrzypek2020}). There are no studies devoted exclusively to \textit{varda} in the history of Swedish; however, for the cognate \textit{verða} in closely related Icelandic we find suggestions that it was used as a futurity marker in the Old and Middle Icelandic period, i.e., 1050--1350 and 1350--1550 respectively (\cite{Halbe1963}: 39--46). For the Swedish material, Svenska Akademiens Ordbok (\citeauthor{SAOB}, The Dictionary of the Swedish Academy) quotes several examples of the use of the verb as a future tense marker, with either infinitive (\ref{ex:skrzypek:4}) or present participle (\ref{ex:skrzypek:5}) complements, some as recent as early 1900s, both constructions are marked as obsolete in the dictionary.

\ea \label{ex:skrzypek:4}
    (\citeauthor{SAOB}, example dated 1595)\\
    \gll At gud alzmectig warder at senda sit harda straff öffuer oss alla.\\
    that God almighty varda\textsc{.prs} to send his.\textsc{refl} severe punishment over us all\\
    \glt ‘That God almighty will send his severe punishment over us all.’
\z

\ea \label{ex:skrzypek:5}
    (\citeauthor{SAOB}, example dated 1928)\\
    \gll En hel del annat som du snart varder skådande.\\
    a whole bit other which you soon become.\textsc{prs} seeing\\
    \glt ‘Other things, which you will soon be seeing.’
\z


The verb \textit{varda} is taken to signify inchoative or ingressive meaning. Some authors, notably \citet{Toyota2009} and \citet{Propst2001}, narrow it down to a mutative meaning – not a general beginning of a new state, but rather a change of state from one to another (see also \cite{Denison1993}:426) or coming into being. Some of the early examples of \textit{varda} illustrate this narrower meaning, as shown in \xref{ex:skrzypek:6}. 

\ea \label{ex:skrzypek:6}
    (Pent, 1300)\\
    \gll Swa lot han oc ewam wardha.\\
    so let he also Eve become\\
    \glt ‘So let he [also] Eve come into being.’
\z

Otherwise, \textit{varda} has been studied mainly as a passive auxiliary (\cite{Markey1969}, \cite{Elmevik1970}, \cite{Skrzypek2020}, \cite{Skrzypek2024}), a function in which it also appears in the oldest extant sources. A lot of research has centered around the question of the loanword \textit{bliva} ‘remain’, which was borrowed into Swedish at some time in the late 1200s (the oldest examples are found in the early 1300s). The verb gradually developed a new meaning ‘become’ alongside the original ‘remain’ and replaced \textit{varda} in all constructions, including the passive one (see \ref{ex:skrzypek:7} and \ref{ex:skrzypek:8}). 

\ea \label{ex:skrzypek:7}
    (Bur, ca. 1330)\\
    \gll þa varþ han takin mz myclom heþar i efeso staþ\\
    then became he taken with much glory in Ephesus town\\
    \glt ‘He was then taken with much glory to the town of Ephesus.’
\z

\ea \label{ex:skrzypek:8}
    (Brahe, 1585)\\
    \gll Anders Person bleff strax greppen på sin gård Ranchyttan\\
    Anders Person became immediately captured on his farm Ranchyttan\\
    \glt ‘Anders Person was immediately captured on his farm Ranchyttan.’
\z

As a result, the usage of \textit{varda} declined, and the verb lives on only in dialects, but is not part of the standard idiom. 

The aim of the present paper is to discuss the history of the verb \textit{varda} in the Swedish prose from 1300 to 1750, in particular whether it could be considered to be a marker of future reference, in combination with non-finite forms of other verbs. The time frame of the study coincides with the major developments of the passive construction and with the documented presence of \textit{varda} in all prosaic genres. We aim to establish whether \textit{varda} could be found in futurate constructions and we will examine the circumstances in which a potential grammaticalization into a future tense marker was halted. The term construction is used in the Construction Grammar (CxG) sense, i.e., we take constructions to be emergent clusters of memory traces, aligned within our conceptual space on the basis of shared form, function, and contextual dimensions (\cite{Goldberg2019}: 7).

The paper is structured as follows: in Section \ref{sec:skrzypek:2} we give a brief account of diachronic studies of future constructions in Germanic languages, with a focus on the German \textit{werden}. Section \ref{sec:skrzypek:3} includes information about sources, methods and tools which were employed to study the development of \textit{varda}. Section 4 gives the general results in terms of frequencies and construction types connected with \textit{varda} in the history of Swedish. In Section \ref{sec:skrzypek:5} we discuss the material in more detail and attempt a qualitative analysis of the examples. Finally, Section \ref{sec:skrzypek:6} concludes the study with a discussion of the reasons for why the Swedish \textit{varda} construction never grammaticalized into an expression of future tense. 

\section{Future tense in a diachronic perspective} \label{sec:skrzypek:2}
Most attempts to define future tense refer to the classical study by \citet[287--298]{Reichenbach1947} and posit that any definition of temporal relations must take into account three locations in time that are positioned relative to each other on a straight line. They are: the moment of speech (S), the moment of an event (E), and a point of reference (R). The latter signifies a vantage point from which an event is viewed, which need not coincide with S. R can be mentioned explicitly, e.g., ‘When we arrived at the station, the train had already left’, or it can be understood from the context. It is possible to define each tense uniquely with reference to these three locations in time, e.g., the present tense is the coincidence of E, R, and S on a single point of the time line.

Although very elegant from a logical and formal point of view, \citeauthor{Reichenbach1947}’s model has some shortcomings which become apparent in historical linguistic studies. Most importantly, it fails to capture constructional polysemy, i.e., the cross-linguistic generalizations that future constructions commonly express other meanings or are derived from other constructions, which are still present in language. Furthermore, it does not account for the fact that many languages have more than one future construction, with different distributional restrictions placed on them (see \cite[17--18]{Hilpert2008} for a discussion and \cite{Comrie1981} for critique). This makes it difficult to apply the model in a diachronic study. 

The rise of future tense markers has mainly been approached from a grammaticalization perspective and, in recent years, within the framework of Diachronic Construction Grammar. It is impossible to refer to the entirety of literature on the topic, but to name just a few relatively recent approaches to English modal verb constructions it includes \citet[]{Traugott1989}, \citet[]{Gotti2003}, \citet[]{Aijmer1985}, \citet[]{Visser1963}, \citet[]{Ziegeler2006}. A number of Germanic future tense markers are investigated diachronically in \citet[]{Hilpert2008}, including the Swedish de-venitive \textit{komma att} and the modal \textit{ska}. The development of \textit{varda} is understandably not included, as the grammaticalization, if at all initiated, did not progress to a fully-fledged marker of future tense. 

There are relatively fewer studies which consider other sources of future tense markers. However, for the Swedish \textit{varda} there is a parallel development in German, with the cognate verb \textit{werden} ‘become’, which has been studied extensively (\cite{Diewald1997}, \cite{Kotin2003}, \cite{Hilpert2008}, \cite{Hartmann2021}, \cite{Kraemer2005}, \cite{Thieroff1992}, \cite{Vater1997} to name just a few). As there are unequivocal similarities between the two verbs, we will briefly recapitulate the developments in German, which will be of relevance for the present study.

German future tense construction with \textit{werden} has developed from an inchoative or ingressive aspect marker (\cite{Diewald1997}: 121). In contrast to the aspectual futures discussed in \citet[]{BybeePerkinsPagliuca1994}, the German \textit{werden}-construction is periphrastic rather than morphological (or zero) and it is a fully grammaticalized future tense marker. The verb itself can also be a copula, which conveys the meaning of an incipient change of state and takes nominal and adjectival complements, e.g. \textit{Peter wird Lehrer} ‘Peter is training to be a teacher.’ or \textit{Peter wird unruhig}. ‘Peter is getting nervous.’ (German examples after \cite[132]{Hilpert2008}). It can be also used as a passive auxiliary, with a past participle, e.g. \textit{Peter wird beobachtet}. ‘Peter is being observed.’ (\cite{Hilpert2008}: 132). We find parallel examples in the Swedish material, here given in the adapted version, \textit{Han vard kung }‘He became a king.’, \textit{Han vard vred} ‘He turned furious.’ and \textit{Han vard gripen} ‘He was captured.’ The similarities between German \textit{werden} and Swedish \textit{varda} (and its later replacement \textit{bliva}) are in this respect striking, as \textit{varda} is found in exactly the same array of functions and has also grammaticalized into a passive auxiliary when accompanied by a passive past participle. 

There is, however, one more construction in Modern German, in which \textit{werden} is found, this time with an infinitive complement, e.g., \textit{Peter wird singen}. ‘Peter will sing’. There is some controversy in the interpretation of the last sentence, which is ambiguous between the temporal interpretation ‘Peter will sing (at some given point in the future)’ and the epistemic interpretation ‘Peter probably sings (right now or generally)’\footnote{Note that the English future construction with \textit{will} may also be used in this way: \textit{We keep telling him to stop but he will do it} can be interpreted as future but also as epistemic, that is just how he is, he does and will keep doing this whatever we say.}. The ambiguity of the construction is the source of much controversy in the account of its development into the future tense construction, dating back at least to \citet[]{Saltveit1962}. \citet[175]{Saltveit1962} argues that perfective verbs such as \textit{kommen} ‘come’, tend to receive a temporal interpretation, because of the implied reference to the end point of an event. Stative verbs such as \textit{sein} ‘be’ on the other hand give the whole construction an epistemic flavour. In this account, Saltveit is followed by a number of authors, although some, notably \citet[196]{Leiss1992}, posit the inverse relation\footnote{It is impossible to give even a most succinct overview of the vast body of research devoted to the grammaticalization of \textit{werden} as future tense marker in the present paper. The reader is directed to \citet[]{Hilpert2008} and \citet{Osterkamp2013} for an exhaustive account of current research as well as novel proposals and to \citet[]{Hartmann2021} for the most recent advances into the formation of the future tense construction. }. In his corpus analysis of the rise of the future construction with \textit{werden}, \citet[133--156]{Hilpert2008} tests whether perfective and continuative verbs are similarly attracted to the construction throughout different periods in time, in order to establish which of the uses, the modal or the temporal one, is the primary one, from which the other is derived. \citet[]{Hilpert2008} shows through a diachronic distinctive collexeme analysis that there is evidence against the account that future meaning developed out of the meaning of intentionality, as suggested in \citet[]{Heine1995}. Further, there is evidence for the primacy of temporal meaning over the epistemic modal one \citep[155]{Hilpert2008}. 

In the diachronic accounts of the rise of \textit{werden}-future it has been further observed that there are in fact two constructions which are in competition, \textit{werden} with an infinitive complement and \textit{werden} with a present participle complement. Some accounts give one primacy over the other, but in the most recent proposal by \citet[]{Hartmann2021} the two constructions are shown to first mutually attract each other and finally merge rather than develop from each other, a process facilitated by the collapse of the infinitive and the present participle into one form. The development entails an expansion in terms of verbs which can occupy the infinitive or present participle slot in the constructions, so that the group of verbs which can be used becomes more heterogeneous in terms of \textit{aktionsart}, beginning with verbs denoting states and activities and gradually encompassing verbs expressing accomplishments and achievements \citep[]{Hartmann2021}. 

As exemplified in \xref{ex:skrzypek:4} and \xref{ex:skrzypek:5}, for convenience repeated as \xref{ex:skrzypek:9} and \xref{ex:skrzypek:10}, \textit{varda} in Swedish could take both infinitives and present participles as complements; there are, however, no corresponding processes of reduction and the two forms are distinct throughout the whole documented history of the Swedish language. 

\ea \label{ex:skrzypek:9}
    (\citeauthor{SAOB} \textit{varda}, example dated 1595)\\
    \gll At gud alzmectig warder at senda sit harda straff öffuer oss alla.\\
    that God almighty varda.\textsc{prs} to send his.\textsc{refl} severe punishment over us all\\
    \glt ‘That God almighty will send his severe punishment over us all.’
\z

\ea \label{ex:skrzypek:10}
    (\citeauthor{SAOB} \textit{varda}, example dated 1928)\\
    \gll En hel del annat som du snart varder skådande.\\
    a whole bit other which you soon become seeing\\
    \glt ‘Other things, which you will soon be seeing.’
\z


No similar convergence could therefore have taken place in Swedish. However, we can study the development in terms of both general frequency of the constructions with the infinitive and the present participle, as well as in terms of the verbs occupying the open slots in the construction (infinitives and present participles). 

The account of previous studies of the diachrony of \textit{werden} is relevant to the subject of the present paper, as there are similarities, both in etymology and use, between the two verbs, which may be reflected in similar developmental paths.

\section{Methods, sources and tools} \label{sec:skrzypek:3}

For the present study we excerpted a total of fifteen Swedish texts written down between 1300 and 1750. We chose prose texts, with both religious and secular prose included. For the sake of the present analysis, the texts were divided into five sub-periods: 1300--1350, 1350--1450, 1450--1550, 1550--1650, and 1650--1750, with three texts in each period. The division of the corpus roughly corresponds to the conventional periodisation of the Swedish language, which is illustrated in Table \ref{tab:skrzypek:1} below. The acronyms used in referring to the texts are explained in the Sources section.

\begin{table}
\caption{The corpus and periodisation of Swedish}
\label{tab:skrzypek:1}
\scriptsize
\begin{tabularx}{\textwidth}{X X X X p{3cm}}
  \lsptoprule
  & {Conventional nomenclature} & {Conventional periodisation} & {Periods in the present study} & {Texts included in the corpus} \\
  \midrule
  \multirow{2}{*}{Old Swedish} & Early & 1225–1350 & 1300–1350 & Bur, Pent, KS \\
                               & Late  & 1350–1526 & 1350–1450 & KM, ST, SVM \\
                               &       &           & 1450–1550 & Linc, Luc, Troj \\
  \midrule
  \multirow{2}{*}{Modern Swedish} & Early & 1526–1732 & 1550–1650 & Brahe, Gyll, Swart \\
                                  & Late  & 1732–     & 1650–1750 & Hiarne, Högström, Spegel \\
  \lspbottomrule
\end{tabularx}
\end{table}

Details of each can be found in the Sources section at the end of this paper. Digitalised versions of the Swedish texts are freely available from \textit{Fornsvenska textbanken }(\href{https://project2.sol.lu.se/fornsvenska/}{https://project2.sol.lu.se/fornsvenska}). In some cases the paper editions of the texts were double-checked to confirm that the digitalised versions do not contain errors. For each of the hundred-year periods, approximately 1800 VPs were annotated, evenly distributed across three texts (per period). In some cases there were slightly more, so that complete portions of text could be analysed. Table \ref{tab:skrzypek:2} shows the size of the corpus and the number of annotated VPs. Over 9000 VPs were tagged with information about (morphological) tense, mood, diathesis, subject animacy, position in sentence and type of clause (main vs subordinate)\footnote{The research presented here is part of a larger project centering around the rise of the periphrastic passive in Swedish and Danish, thus the paper does not utilize all annotated information. The tool used is DiaPass, created specifically for the project.}. In the present paper, this data is used mainly to illustrate the relative frequency of \textit{varda} and all verbs, as well as the relative frequency of \textit{varda} and \textit{bliva}. As \textit{varda} is not high in frequency in proportion to all verbs, we have gleaned only 180 examples from the entire annotated material. This gives us some idea of the verb’s constructional potential (i.e., what constructions it can be found in), but is not enough to follow its development. Therefore this study was completed with manual excerption of the 15 texts in their unabridged form. They were searched for all occurrences of \textit{varda} in their entirety\footnote{All spelling variants and tense forms of \textit{varda} were included in the search; it is however possible that a handful of examples were missed in the search, which should not affect the overall results.}. The search returned 1422 examples, which will be analysed qualitatively. 

\largerpage[1]
\newpage

\begin{table}
\caption{The material}
\label{tab:skrzypek:2}
\scriptsize
\begin{tabularx}{\textwidth}{lCCC c r}
  \lsptoprule
  {Text} & {Entire text no of words} & {No of instances of \textit{varda}} & {\textit{varda} per 1000 words} & {Date (assumed)} & {Period} \\
  \midrule
  Pent & 140760 & 356 & 2.53 & 1300 & I (1300--1350) \\
  Bur & 191808 & 33 & 0.17 & 1330 & I (1300--1350) \\
  KS & 21600 & 74 & 3.43 & 1330 & I (1300--1350) \\
  KM & 10930 & 29 & 2.65 & 1380 & II (1350--1450) \\
  SVM & 16240 & 40 & 2.46 & 1420 & II (1350--1450) \\
  ST & 145479 & 497 & 3.42 & 1420 & II (1350--1450) \\
  Lucidarius & 28175 & 151 & 5.36 & 1487 & III (1450--1550) \\
  Linc & 9824 & 37 & 3.77 & 1520 & III (1450--1550) \\
  Troj & 44380 & 8 & 0.18 & 1529 & III (1450--1550) \\
  Swart & 51940 & 61 & 1.17 & 1560 & IV (1550--1650) \\
  Brahe & 26380 & 4 & 0.15 & 1585 & IV (1550--1650) \\
  Gyllenhielm & 53010 & 113 & 2.13 & 1640 & IV (1550--1650) \\
  Hiarne & 11030 & 7 & 0.63 & 1660 & V (1650--1750) \\
  Spegel & 32720 & 12 & 0.37 & 1680 & V (1650--1750) \\
  Hogstrom & 17990 & 0 & 0.00 & 1747 & V (1650--1750) \\
  \lspbottomrule
\end{tabularx}
\end{table}

The number of hits for \textit{varda} in each text varies not just according to the text’s length (which is trivial, as we expect fewer hits in shorter texts), and not just according to the chronology, although we see a tendency for \textit{varda} to gradually disappear throughout the periods studied. It seems also that, especially after 1450, the use of \textit{varda} was very much dependent on the scribe’s preferences (or rather dialect). Thus Brahe (1585) and Hiarne (1660) show very low values, while Gyllenhielm (1640) uses the verb often. We return to these differences and their likely origins in Section \ref{sec:skrzypek:4}. 

It is notoriously difficult to establish the meanings of constructions in historical texts. There is the risk of imposing modern reading on the material, as well as a lack of clear diagnostics in the context, which would clarify the meaning of a given structure. This problem becomes paramount when incipient structures are considered, as we expect them to be ambiguous and constitute bridging contexts (in the sense of \cite{Heine2002}), and an unambiguous analysis of such examples may not be possible. 

There seems to exist only one unequivocal diagnostic of future reference, which is the presence of temporal adverbials indicating a future location in time. Disregarding the construction type (be it present tense, modal verb, verb of motion, etc.), if such an adverbial is present in the clause, the entire clause can be taken to have futurate reference. 

A problem which nevertheless presents itself is the distinction between the immediate and the remote future. While a temporal adverbial indicating remote future, e.g. \textit{on the day of judgement}, may give an unambiguous future reading (see \ref{ex:skrzypek:11}), it is less clear whether adverbials indicating immediate future, e.g. \textit{immediately}, result in the same reading (compare with \ref{ex:skrzypek:12}). 

\ea \label{ex:skrzypek:11}
    (Pent, ca. 1300)\\
    \gll Nw sighir war hærra innan læstinne at wi vardhom like ænglom æptir domadagh.\\
    now says our lord within script.\textsc{def.dat} that we become.\textsc{pl} like angel.\textsc{pl.dat} after doomsday\\
    \glt ‘Now our Lord says in the script that we will be like angels after the Day of Judgement.’
\z

\ea \label{ex:skrzypek:12}
    (Brahe, 1585)\\
    \gll efter theras fader war slagin i krigh, som strax omtaladt warder.\\
    after their father was killed in war, which immediately discussed becomes\\
    \glt ‘After their father was killed in war, which will immediately be discussed.’
\z

It may also be said that the futurate interpretation of \xref{ex:skrzypek:12} arises from the present tense and is mainly aspectual.  

In this section we have presented the material used for the present study and discussed some aspects of approaching the data. The following sections are devoted to presenting an overview of the quantitative results (Section \ref{sec:skrzypek:4}) and a more detailed and mainly qualitative analysis (Section \ref{sec:skrzypek:5}). 

\section{\textit{varda} in Swedish -- an overview} \label{sec:skrzypek:4}

In this section we present the use of \textit{varda} in Swedish between 1300 and 1750. As shown in \tabref{tab:skrzypek:3}, \textit{varda} is not a common verb in the OSw material. \tabref{tab:skrzypek:3} illustrates the frequencies of both \textit{varda} and its competition, the loanword \textit{bliva}, in each of the periods studied. As we can see, the frequency of \textit{varda} is relatively stable between 1300 and 1550 and falls significantly after 1650, while \textit{bliva} gains significantly in frequency after 1550. 

\begin{table}
\caption{Frequencies of \textit{varda} and \textit{bliva} in the corpus}
\label{tab:skrzypek:3}
\small
\begin{tabularx}{\textwidth}{lCCCCC}
  \lsptoprule
  {Period} & {Entire text no of words} & {\textit{varda}} & {\textit{varda} per 1000 words} & {\textit{bliva}} & {\textit{bliva} per 1000 words} \\
  \midrule
  1300--1350 & 354168 & 463 & 1.31 & 7 & 0.02 \\
  1350--1450 & 172649 & 566 & 3.28 & 211 & 1.22 \\
  1450--1550 & 82379  & 196 & 2.38 & 166 & 2.02 \\
  1550--1650 & 131330 & 178 & 1.36 & 926 & 7.05 \\
  1650--1750 & 61740  & 19  & 0.31 & 328 & 5.31 \\
  \midrule
  \textbf{Total} & 802266 & 1422 &  & 1638 &  \\
  \lspbottomrule
\end{tabularx}
\end{table}

The verb \textit{varda} could be used in a wide array of constructions, taking both verbal (infinitive INF, past participle PTCP or present participle PPart) and nominal (NP, AdjP) complements. Representative examples are quoted below (\ref{ex:skrzypek:13}--\ref{ex:skrzypek:17}).

\largerpage[-1.5]

\ea \label{ex:skrzypek:13}
    [SBJ \textit{varda} INF] (KM, 1400)\\
    \gll Nw wardhom wi tala om them som fore waro stadde.\\
    now become.\textsc{1pl} we talk.\textsc{inf} about them who before were met\\
    \glt ‘We shall now tell about those who were met before.’
\z

\ea \label{ex:skrzypek:14}
    [SBJ \textit{varda} PTCP] (Bur, ca. 1330)\\
    \gll þa varþ han takin mz myclom heþar i efeso staþ.\\
    then became he taken with much glory in Ephesus town\\
    \glt ‘He was then taken with much glory to the town of Ephesus.’
\z

\ea \label{ex:skrzypek:15}
    [SBJ \textit{varda} PPart] (SVM, 1420)\\
    \gll Aff hwilko ther riddare-n wardh syrghiande.\\
    of which there knight-\textsc{def} became grieving\\
    \glt ‘What made the knight grieve.’
\z

\ea \label{ex:skrzypek:16}
    [SBJ \textit{varda} NP] (Bur, 1330)\\
    \gll tel han varþ røuare ok røuar-a mæstare.\\
    until he became robber and robber-\textsc{pl.gen} master\\
    \glt ‘Until he turned into a robber and master of the robbers.’
\z

\ea \label{ex:skrzypek:17}
    [SBJ \textit{varda} AdjP] (SVM, 1420)\\
    \gll Tha wardh swin-ith spakare.\\
    then became pig-\textsc{def} clever.\textsc{comp}\\
    \glt ‘Then the pig became more clever.’
\z


The loanword \textit{bliva} is found in the same constructions, although it takes about 200 years since its introduction into Swedish for it to be used with all types of complements (\cite{Skrzypek2024}). Throughout the entire period studied, \textit{bliva} can also be used with locative adverbials. This use is clearly connected with the verb’s original meaning ‘remain’, see \xref{ex:skrzypek:18}. 

\ea \label{ex:skrzypek:18}
    [SBJ \textit{varda} PTCP] (Hiarne, 1665)\\
    \gll Och alltså bleef Celadon uthi Epheso.\\
    and thus stayed Celadon in Ephesus\\
    \glt ‘And so Celadon stayed in Ephesus.’
\z

No corresponding examples of \textit{varda} with locatives as complements were found in the material.

There are also less common uses of the verb \textit{varda}, which are mainly found in the oldest material. They include an ‘absolute’ use, without any complements, when the verb can be taken to mean ‘come into existence’, see (\ref{ex:skrzypek:6}) above and (\ref{ex:skrzypek:19}) below. 


\ea \label{ex:skrzypek:19}
    [SBJ \textit{varda} INF] (ST, 1420)\\
    \gll tha wardh swa mykyt thrang a gatumen / at the ey kundo fram koma.\\
    then became so much throng on street.\textsc{dat.pl.def} {} that they not could forward come\\
    \glt ‘Then (suddenly) there was such a throng on the streets that they could not come out.’
\z

This construction is found mainly in Pent, the oldest text in the sample; some instances were identified in SVM and Luc as well, but none after 1550 in the corpus.

In Table \ref{tab:skrzypek:4} we show the distribution of \textit{varda} across the constructions which were most common in the corpus, divided by periods. We give the absolute numbers of occurrences, as the percentage values are too low. 

\begin{table}
\caption{Occurrences of \textit{varda} in the major constructions }
\label{tab:skrzypek:4}
\begin{tabularx}{\textwidth}{lCCCCCC}
  \lsptoprule
 & {[SBJ \textit{varda} NP]} & {[SBJ \textit{varda} AdjP]} & {[SBJ \textit{varda} PTCP]} & {[SBJ \textit{varda} PPart]} & {[SBJ \textit{varda} INF]} & {Total} \\
  \midrule
  1300--1350 & 78 & 145 & 94 & 9 & 22 & 348 \\
  1350--1450 & 66 & 199 & 195 & 17 & 11 & 488 \\
  1450--1550 & 16 & 44  & 95 & 6 & 0  & 161 \\
  1550--1650 & 9  & 17  & 122 & 15 & 3  & 166 \\
  1650--1750 & 1  & 2   & 12  & 3  & 0  & 18  \\
  \midrule
  \textbf{Total} & 170 & 407 & 518 & 50 & 36 & 1181 \\
  \lspbottomrule
\end{tabularx}
\end{table}

The quantitative overview shows that \textit{varda} can appear in a number of constructions throughout the period studied, including constructions which are strong candidates for expressions of future tense, especially with regard to similar developments taking place in the closely related German and with the cognate \textit{werden} at the centre (e.g. \cite{Hartmann2021}). We may note, however, that by 1450 \textit{varda} is mainly found with past participle complements and the entire construction is passive (much like in German). Development beyond 1550 involves mainly the suppression of \textit{varda} in favour of \textit{bliva }(\cite{Skrzypek2024}). 

In the following section we will present a qualitative analysis of the potential future constructions with \textit{varda}. 

\section{\textit{varda} in Swedish -- a detailed analysis} \label{sec:skrzypek:5}

As we have discussed in Section \ref{sec:skrzypek:2}, it is challenging to discern future reference in historical material and it is only seldom that the context provides enough cues to be able to do so with any confidence. Linguistic meaning is notoriusly difficult to quantify and measure objectively, but in this section we will attempt to discuss in more detail the material collected through manual excerption, sorted by construction type, including the three constructions that make the best candidates for a future tense construction, i.e., one with the infinitive, one with the present participle and the aspectual future. In Table \ref{tab:skrzypek:5} we show the frequencies of [SBJ \textit{varda} INF] and [SBJ \textit{varda} PPart] across the texts in the sample.

\begin{table}
\caption{Occurrences of \textit{varda} with the infinitive and with the present participle}
\label{tab:skrzypek:5}
\begin{tabularx}{\textwidth}{lCCCC}
  \lsptoprule
  {Text} & {Entire text no of words} & {No of instances of \textit{varda}} & {[SBJ \textit{varda} INF]} & {[SBJ \textit{varda} PPart]} \\
  \midrule
  Pent & 140760 & 356 & 22 & 4 \\
  Bur  & 191808 & 33  & 0  & 1 \\
  KS   & 21600  & 74  & 0  & 2 \\
  KM   & 10930  & 29  & 2  & 0 \\
  SVM  & 16240  & 40  & 0  & 1 \\
  ST   & 145479 & 497 & 9  & 19 \\
  Lucidarius & 28175  & 151 & 0  & 2 \\
  Linc & 9824   & 37  & 0  & 1 \\
  Troj & 44380  & 8   & 0  & 1 \\
  Swart & 51940  & 61  & 0  & 13 \\
  Brahe & 26380  & 4   & 0  & 2 \\
  Gyllenhielm & 53010  & 113 & 3  & 1 \\
  Hiarne & 11030  & 7   & 0  & 3 \\
  Spegel & 32720  & 12  & 0  & 0 \\
  Hogstrom & 17990  & 0   & 0  & 0 \\
  \lspbottomrule
\end{tabularx}
\end{table}

The most likely candidate for the future expression is the construction with infinitive as a complement of \textit{varda}, [SBJ \textit{varda} INF]. It is exemplified in (\ref{ex:skrzypek:20}) below, see also (\ref{ex:skrzypek:13}) in Section \ref{sec:skrzypek:4}. 

\ea \label{ex:skrzypek:20}
    (Pent, 1300)\\
    \gll Nw varda the gøra som han sagde.\\
    now become they do as he said\\
    \glt ‘Now they will do as he said.’
\z

There are only 36 similar examples in the entire corpus (see Table \ref{tab:skrzypek:5}), the majority in the oldest text in the sample, Pent. This, however, is not the only text where we find \textit{varda} with the infinitive. There are some examples in KM (1400), see (\ref{ex:skrzypek:21}), ST, see (\ref{ex:skrzypek:22}) and, somewhat surprisingly, as late as in 1640, in Gyll, see (\ref{ex:skrzypek:23}). In this respect Gyll is completely isolated, and most likely the use of \textit{varda} there is strongly influenced by German.\footnote{In his study of word order in Swedish, \citet[]{MagnussonPetzell2014} shows some examples from Gyll which strongly suggest German influence on the syntactic level visible in the text. It is not unlikely that such influence could also include constructional interference with the verb \textit{varda}, due to its similiarity to \textit{werden}.}

\ea \label{ex:skrzypek:21}
    (KM, 1400)\\
    \gll Nw wardhom wi tala om them som fore waro stadde.\\
    now become we speak about them who before were set\\
    \glt ‘We will now speak of those who were living before.’
\z

\ea \label{ex:skrzypek:22}
    (ST, 1420)\\
    \gll swa længe gudh ær j hymerike tha wardher iak her swa brinna.\\
    so long god is in heaven then become I here so burn\\
    \glt ‘As long as God is in heaven, I will burn here so.’
\z

\ea \label{ex:skrzypek:23}
    (Gyll, 1640)\\
    \gll som historien uthvijsa varder.\\
    as history show will\\
    \glt ‘As history will show.’
\z

In KM, the construction is only found sporadically and seems to be restricted lexically, so that only infinitives of \textit{verba dicendi} can appear there. It seems therefore that the construction was present in the oldest material\footnote{\citeauthor{SAOB} quotes more recent examples of \textit{varda} with infinitive, the youngest are from the 1500s (see also \xref{ex:skrzypek:4}.}, but rather uncommon, and its later occurrences are a result of the influence of German. 

We have remarked before that the borrowed \textit{bliva}, originally found only with locative adverbials (see \ref{ex:skrzypek:18}), gradually developed the polysemous meaning ‘remain/become’ and started appearing with the same complements as \textit{varda}, i.e., with NPs, AdjPs and participles. However, we have not found a single example of \textit{bliva} with the infinitive, which suggests that by the time the polysemy arose, the construction with the infinitive was no longer found in the material and the new verb could not be used with infinitives. 

Even though rarely found, the construction with the infinitive comes closest to what might be regarded an (incipient) expression of futurity. Discussing the grammaticalization of the future tense, \citet[]{BybeePerkinsPagliuca1994} state the following:

\begin{quote}
    We hypothesize that all futures go through a stage of functioning to express the intention, first of the speaker, and later of the agent of the main verb (Bybee and Pagliuca \citeyear{bybeeEvolutionFutureMeaning1987}, Bybee \citeyear{bybeeSemanticSubstanceVs1988}). The meanings that can feed the future path must be meanings that appropriately function in statements that imply an intention on the part of the speaker. (\cite{BybeePerkinsPagliuca1994}: 254)
\end{quote}

It is not entirely clear whether there is any intentionality on the part of the subjects of the sentences we find where \textit{varda} is complemented by an infinitive, with the exception of sentences with \textit{verba dicendi }and first person subjects (as in \ref{ex:skrzypek:13}). These contexts, however, are all of the subject being identical with the speaker. The structure can be seen as a bridging context, allowing for a broader interpretation of volition and intention as not restricted to the speaker but rather as characterizing the subject. In use with other verbs, the subject may not be identical with the speaker, yet retain intentionality (see \citet{Heine1995} for a similar explanation of the development of the German \textit{werden}, based on synchronic considerations and \citet{Hilpert2008} for a diachronic corpus account).

In sentences where the subject and the speaker are not identical (see \ref{ex:skrzypek:20}), what we find may equally well be the original, lexical meaning of \textit{varda}, i.e., ‘to begin’. Translated into English, \textit{varda} in these examples can mean ‘will’, but it can also mean ‘begin’. It seems that to consider them futurate, we would impose an anachronistic reading. However, in examples similar to \xref{ex:skrzypek:22}, ‘as long as God is in heaven, I will burn here’, the wording strongly implies that the burning has already started (and is not about to begin) and the speaker expects it to continue for a time yet, rendering the construction futurate. Such examples are confined to the earliest texts in the corpus. However, they seem to have been productive (in the sense of \citet{Barddal2008}, based on the concept of extensibility, i.e., being used with new verbs) and not merely calqued into translated texts from the original ones, as we find examples of futurate \textit{varda} with infinitive also in the provincial legal texts from late 1200s.

\xref{ex:skrzypek:23} on the other hand comes from a relatively late text, from 1640. It is therefore striking. The only explanation of its existence (and two more similar examples can be found in the same text) that we can offer is that it is the author’s proficiency in German (see also \cite{MagnussonPetzell2014}) that allows him to use this construction. This is strengthened by the fact that the examples in Gyll are isolated, and similar constructions were not found in other contemporary texts. 

Summing up, there are too few examples of [SBJ \textit{varda} INF] to categorically state that what we see is incipient grammaticalization, but enough to demonstrate that the construction was present in the language use, albeit under some constraints, at least in the early 1300s, and that it was acceptable enough to be used in the translation of ST in the early 1400s. Its occurrence in Gyll in 1640 is very difficult to explain unless we assume German influence. The most diverse examples (in terms of verbal complements) are found in Pent, the oldest text in the sample. Many authors who have studied its linguistic form have noted its archaic nature, in particular \citet{Hesselman1927}\footnote{Då man läser något parti av MBI, särskilt i de berättande styckena, t.ex. delar av Genesis, får man rätt snart ett högst tilltalande intryck av en mycket gammaldags och kärnfull svenska, en språklig äkthet och på samma gång rikedom, som man eljest icke lätt finner motstycken till i vår gamla litteratur. \citep[21]{Hesselman1927}}, see also \citet{Thorell1959}. Furthermore, Pent is generally considered free from foreign influence \citep{Wollin2001}. This suggests that the use of \textit{varda} in this material represents an authentic construction. 

Another potential candidate for a future tense construction is one with the present participle in the complement position, [SBJ \textit{varda} PPart]. It is exemplified in \xref{ex:skrzypek:24} below.

When one reads an excerpt from MBI, especially in the narrative passages, e.g. parts of Genesis, one gets a highly pleasing impression of a very old-fashioned and pithy Swedish, a linguistic authenticity and at the same richness which one otherwise does not easily find parallels to in our old literature. 

\ea \label{ex:skrzypek:24}
    (Petri, 1530)\\
    \gll Allom them som thenna Cröneke läsande eller hörande warda, önskar jach Olauus Petri, predikere i Stocholm, saligheetennes kundtskap.\\
    all.\textsc{dat} them who this chronicle reading or hearing varda wish I Olaus Petri preacher in Stockholm bliss.\textsc{gen} wisdom\\
    \glt ‘All those who are/will be reading or hearing this chronicle, I, Olaus Petri, preacher in Stockholm, wish the wisdom of bliss.’
\z

This construction is reminiscent of the German construction with \textit{werden} and the present participle, which under a long period of time competed with the construction with the infinitive, for the status of a futurity marker (\citet{Hartmann2021} and references therein). Most likely due to a blurring of present participle and infinitive forms (\textit{lachend} vs. \textit{lachen }‘laughing’ vs. ‘(to) laugh’, see \cite{Hartmann2021}: 367) the construction was attracted to the one with the infinitive and the two converged, resulting in [SBJ \textit{werden} INF] as the only marker of futurity in Modern German (\cite{Hartmann2021}). 

In [SBJ \textit{varda} PPart] the mutative, inchoative or ingressive meaning is apparent. In each of the examples (\ref{ex:skrzypek:25}--\ref{ex:skrzypek:27}) below, we see an event presented by means of [SBJ \textit{varda} PPart] construction, which is shown to be the result of an earlier event (suggested in the context by \textit{aff hwilko} ‘by which’, \textit{nar han saat} ‘when/after he sat down’ and by \textit{efter thett att} ‘after that’). 

\ea \label{ex:skrzypek:25}
    (SVM, 1420)\\
    \gll Aff hwilko ther riddaren wardh syrghiande.\\
    of which there knight-\textsc{def} became grieving\\
    \glt ‘What made the knight grieve.’
\z

\ea \label{ex:skrzypek:26}
    (Linc, 1520)\\
    \gll tha böriadhe hälga iomffrun til at mykyt sokkas oc gratandhe wardha.\\
    then began saint virgin.\textsc{def} to to much sigh and crying become\\
    \glt ‘And the saint virgin started to sigh and cry.’
\z

\ea \label{ex:skrzypek:27}
    (Brahe, 1585)\\
    \gll Theras fader war slagin i krigh, som strax omtaladt warder.\\
    their father was killed in war which immediately discussed become\\
    \glt ‘Their father was killed in the war, which will now be discussed.’
\z

We can see that the [SBJ \textit{varda} PPart] construction has the potential to signal futurity, or at least, to show one event as taking place as a result of (and therefore also later than) another, earlier event. However, the [SBJ \textit{varda} PPart] construction does not show any signs of grammaticalizing beyond the resultative meaning (as witnessed by \ref{ex:skrzypek:27} from late 1500s) and gradually disappears as \textit{varda} gives way to \textit{bliva}. 

Finally, we consider the constructions in which \textit{varda} is found in the present tense, [SBJ \textit{varda}\textsubscript{pres}], irrespective of the complement, which means we include also instances of the periphrastic passive construction [SBJ \textit{varda} PTCP], as exemplified in \xref{ex:skrzypek:28}. 


\ea \label{ex:skrzypek:28}
    (SVM, 1430)\\
    \gll Waktin fore thy idher thz i wardhin ey swik
in aff gire.\\
    guard for it.\textsc{dat} you that you varda not disappointed of greed\\
    \glt ‘Be on your guard so that you will not be let down by greed.’
\z

A first observation is that there is a significant difference between time before 1450 and after this date. We have noted before that after 1450, \textit{varda} appears mainly with past participles and these constructions are instances of passive voice. Interestingly enough, this change coincides with a change in the distribution of \textit{varda} amongst different tenses. The numbers, based on the annotated material of ca. 9000 VPs, are given in \tabref{tab:skrzypek:6}. Examples with modal verbs were excluded from the analysis, as in these it is the modal verb that is the locus of temporal information.

\begin{table}
\caption{Distribution of \textit{varda} across tenses}
\label{tab:skrzypek:6}
\footnotesize
\begin{tabularx}{\textwidth}{lCCCC}
  \lsptoprule
  {Period} & {Present} & {Past} & {Perfect} & {Past perfect} \\
  \midrule
  1300--1350 & 38.10\% & 52.38\% & 0\%    & 0\%    \\
  1350--1450 & 31.11\% & 66.67\% & 0\%    & 0\%    \\
  1450--1550 & 6.52\%  & 76.09\% & 0\%    & 15.22\% \\
  1550--1650 & 7.14\%  & 85.71\% & 0\%    & 7.14\%  \\
  1650--1750 & 0\%     & 93.75\% & 6.25\% & 0\%    \\
  \lspbottomrule
\end{tabularx}
\end{table}

\largerpage[2]

The sharp fall in numbers after 1450 is partly to be seen as a result of there generally being very few clauses in the present tense in the material in the period 1550--1650. However, in 1450--1550 and 1650--1750,  the percentage of clauses in the present tense is much higher than the numbers for \textit{varda} would suggest. For the sake of clarity, in \tabref{tab:skrzypek:7} we show the results for the present tense first for all annotated VPs (including any finite verb, with the exception of modal verbs), then those with \textit{varda} (repeated from \tabref{tab:skrzypek:6}).

\begin{table}
\caption{The present tense VPs in the corpus}
\label{tab:skrzypek:7}
\begin{tabularx}{0.5\textwidth}{CCC}
  \lsptoprule
  {Present tense} & {All VPs} & {\textit{varda}} \\
  \midrule
  1300--1350 & 47.06\% & 38.10\% \\
  1350--1450 & 25.76\% & 31.11\% \\
  1450--1550 & 31.79\% & 6.52\%  \\
  1550--1650 & 3.87\%  & 7.14\%  \\
  1650--1750 & 18.79\% & 0\%     \\
  \lspbottomrule
\end{tabularx}
\end{table}

Tables \ref{tab:skrzypek:6} and \ref{tab:skrzypek:7} show that the frequency of \textit{varda} in the present tense is on a comparable level to the frequency of all verbs in the present tense until 1450. Then a sharp drop occurs, which cannot be explained by the lower values for the present tense alone (although period 1550--1650 seems to be rather extreme and \textit{varda} reaches higher values than all verbs). What mainly transpires from these data is that \textit{varda} is by no means strongly affiliated with the present tense, apart from the first two periods studied. 

How is the distribution of \textit{varda} in different tenses, in particular between the present and the past tense, relevant to the present study? We would like to argue that the potential futurate meaning of constructions with \textit{varda} is aspectual. Since \textit{varda} denotes a change of state (mutative meaning), or (in many contexts) can be read as ‘begin, come into existence’, in combination with the present tense it will render the whole construction futurate, as the change is usually not instant, but takes some time, shorter or longer. Therefore the speakers are aware of the result of the action (the end-point of the change) being located further from the moment of speech, not being simultaneous with it. The idea of future reference is thus inherent in the meaning of the verb and the future, relative to S (speech time), and becomes fully actualized in the present tense. 

We are thus left with a last diagnostic of possible future reference, the presence of temporal adverbials.

\largerpage[2]

\ea \label{ex:skrzypek:29}
    (Pent, 1300)\\
    \gll thy at sanctus gregorius sighir aff oss at hælge mæn æptir domadagh vardha swa lyusir oc gønom skærir at (…)\\
    it.\textsc{dat} that sanctus Gregorius says of us that holy men after doomsday become so light and through transparent that (…)\\
    \glt ‘Thus Saint Gregory tells us that holy men (will) become so light and transparent after the Day of Judgement that (…)’
\z


We note that the verb \textit{varda} is given in the present tense in these examples. It is very difficult, if not outright impossible, to tease apart the different cues through which the future reference seems to arise. Apart from the verb \textit{varda} itself, and the present tense form, we find adverbials such as \textit{strax} ‘immediately’. Also the adverbial \textit{nw} ‘now’ could be taken to bring to the sentence the reading of immediate future, see \xref{ex:skrzypek:21}. 

In this section we have discussed examples of potentially futurate constructions with \textit{varda}, i.e., constructions with the infinitive or the present participle or the aspectual uses of the verb. The common denominator of the constructions is that the verb’s lexicality is present, i.e, it means ‘become’ or ‘come into being’ or even ‘begin’ and the whole construction is inchoative or ingressive. 

In the following section we bring all the results together to discuss to what extent \textit{varda} could be considered a future tense marker in OSw and why there was no further grammaticalization. 

\section{Discussion} \label{sec:skrzypek:6}
The aim of the present paper was to give a brief history of the verb \textit{varda} in Swedish and, given that some of its uses in the oldest extant texts is reminiscent of nascent future tense constructions, examine the circumstances in which this potential grammaticalization into a future tense marker was halted. 

We have seen examples of \textit{varda} being used as a marker of future tense (with infinitival complements), reminiscent of the development that has taken place in German (of a cognate verb \textit{werden}). In other words, we have some arguments beyond the cross-linguistic comparison and the developments in a closely related language to support the hypothesis that \textit{varda} can be treated as an incipient futurity marker, whose grammaticalization into a fully-fledged exponent of future reference is for some reason halted. The examples are sporadic, and concentrated in the oldest extant texts, with the exception of ST from ca. 1420, which is a translation of a MLG text, heavily influenced by the MLG original and Gyll from 1640, where many other syntactic influences of German can be found. These examples are all of the construction with the infinitival complement, illustrated in (\ref{ex:skrzypek:22}) and (\ref{ex:skrzypek:23}).

It is not possible to discuss \textit{varda} without making a reference to \textit{bliva}, a verb borrowed from MLG in the late 1200s, which gradually takes over all of \textit{varda}’s original domains, with one significant exception. As we have noted, there were no examples found with \textit{bliva} occupying the finite verb slot in the construction in the entire material. It seems that this is the only construction where \textit{varda} could be found in which \textit{bliva} does not suppress it, which suggests that by the time \textit{bliva} developed towards the polysemous meaning ‘remain/become’ and started appearing in constructions with adjectival and nominal complements, the construction with the infinitive was no longer found in the material and the new verb could not be used with infinitives. Interestingly enough, there is contemporary evidence for the rise of a futurate construction with \textit{bliva} and an infinitival complement in Norwegian (see Mikkelsen, this volume). To my knowledge, similar uses of \textit{bliva} are not found in MSw. 

Another candidate for a future tense construction is one with the present participle in the complement position. This construction survives as long as the verb does in the standard language, and its counterpart in MSw is [SBJ \textit{bliva} PPart]. The construction is slightly more common than the one with the infinitive, but its overall usage remains low throughout the entire period studied. In contrast to German, there is no attraction between the constructions, and the present participle and the infinitive are formally distinct in MSw. As a consequence, any future reference which might be inferred from the construction is a result of the verb \textit{varda} used in the present tense rather than the use of the present participle. 

On the other hand, we see \textit{varda} become increasingly associated with the slot in a participial construction (with past rather than present participles), at the expense of its earlier association with other verbal, adjectival or nominal complements (see \tabref{tab:skrzypek:4}). The mutative, or ingressive, meaning is present throughout the period studied, but since the verb is more and more often found with passive participles (which presuppose an agent for the action), the original non-agentive meaning of the verb (‘it came to pass’, ‘it happened’ rather than ‘it was done’) is backgrounded and \textit{varda} can more easily enter into constructions with participles derived from highly agentive verbs, shifting the entire construction towards passive rather than inchoative, ingressive or resultative ones. We may see that there is a parallel here with the grammaticalization of the future tense marking, where there is an increase in intentionality, which is closely intertwined with agentivity. In fact, both can be seen as two sides of the same coin: while agentivity is mainly focused on the past or its present results, intentionality concerns the future, yet in both cases we find that the Agent is in control. However, the orientation of the verb is different: while in the passive voice construction, it is oriented towards the subject that is not the Agent but rather the Patient (and thus does not entail volition or intentionality\footnote{As this paper concerns the (potential) future tense constructions, we cannot go into too much detail concerning the passive voice construction here. However, the periphrastic passive construction [SBJ \textit{bliva} PTCP] in MSw does entail some intentionality on the part of the subject (= the patient), in contrast to a morphological passive construction [SBJ V\textsubscript{pass}], which does not. The reader is directed to \citet{Engdahl1999} and \citet{Laanemets2012} for a more detailed account on the synchronic aspects and to \citet{Skrzypek2024} for diachronic ones.}), in the future tense construction, it is oriented towards the subject that is the Agent and capable of intentional activity.

It is naturally not a given that it was this process alone that halted a potential development towards future marking, as the German \textit{werden} does grammaticalize in both functions. Rather, it seems that the lack of a similar development in Swedish is a result of several factors. Firstly, the nascent future constructions (with the infinitive or the present participle) have a very low usage throughout the periods studied. Secondly, as soon as \textit{bliva} develops the polysemous meaning ‘to remain / to become’, all the developments in \textit{varda} are closely linked to what is happening with \textit{bliva}, and its further development into a passive auxiliary. From being an antonym to \textit{varda} (in the sense of \citet[93]{Lobner1990}), \textit{bliva} gradually converges with it, entering into the same array of constructions in which \textit{varda} was previously found. Significantly, however, not a single example of \textit{bliva} with the infinitive is attested, which could suggest that the construction itself was of a highly limited use. 

Most importantly, however, at no stage does the use of \textit{varda} involve the notion of intentionality, apart from the few examples with the infinitive. Admittedly, some authors have pointed out that intentionality need not be part of the development, e.g. \citet[]{Traugott1978} and \citet[]{Dahl2000}, which may instead proceed through inchoative, ingressive or resultative meanings. These meanings are present in all uses of \textit{varda}. Yet it seems that only when coupled with some measure of intentionality of the subject they can lead to the development of the future reference.

\section{Conclusions} \label{sec:skrzypek:7}
We have discussed the developments and ultimate demise of a construction [SBJ \textit{varda} INF], which was present in the oldest extant Swedish texts, and could be argued to have a futurate reference, but which quickly disappears from the documented language use. At the same time, we see that where \textit{varda} is retained, it can still appear with infinitival complements as late as in the early 1900s (according to \citeauthor{SAOB}). Even though the futural use never becomes part of the standard idiom, it is present in the Swedish language for several centuries. One of the reasons for the incomplete grammaticalization must be sought in the replacement of \textit{varda} by \textit{bliva}. Had it been retained, its original constructional potential, in particular the possibility to take an infinitive complement, could have been strengthened by the renewed contact with German (Middle Low German, as illustrated by the ST and High German, as seen in the 1600s), with its already grammaticalized future tense \textit{werden}-construction. It is of course ‘the land of might-have-been’. Nevertheless, the study of change that did not happen, despite there being a potential for a certain change, is an interesting exercise in language history and the forces which influence language change. Another candidate for such a study would be the Swedish \textit{vilja} ‘will’, which appears with future-like uses in OSw, but then develops into the marker of volition alone (unlike the Danish and Norwegian \textit{vil}). 

We know that \textit{varda} does not grammaticalize into a future tense marker, but our intention was to test the possible differences in context and use which may have led to this divergence of closely related languages, Swedish and German. Naturally, language change which takes place in one language need not translate into identical or even similar developments in another. Yet there are developmental patterns we repeatedly find cross-linguistically, which synchronically take the form of implicational universals and diachronically of paths or clines. Much research is devoted to judging why a change takes place, but it is also relevant to ask why it does not, when apparently there is room for it. 

The conclusion is that to understand language change we need to consider not only the particular construction, which is the locus of change, but also the constructional potential of its components. What happens to a given item outside a given construction is of consequence for language change. Either the divergent development results in a split in the lexeme (separate lexical entries for a given item, originally one word) or the course of the change may be stayed or altered. 

\largerpage[1]

\section*{Abbreviations}
\begin{tabularx}{.5\textwidth}{@{}lQ@{}}
\textsc{1} & first person \\
\textsc{comp} & complementizer \\
\textsc{dat} & dative \\
\textsc{def} & determiner \\
\textsc{gen} & genitive \\
\end{tabularx}%
\begin{tabularx}{.5\textwidth}{@{}lQ@{}}
\textsc{inf} & infinitive \\
\textsc{pl} & plural \\
\textsc{prs} & present \\
\textsc{pass} & passive \\
\textsc{refl} & reflexive \\
\end{tabularx}


\section*{Acknowledgements}
This research was supported by the Polish National Science Center (grant number 2021/43/B/HS2/00048). I would like to thank the editors and the anonymous reviewers for their valuable and constructive comments.


\section*{Sources}

\begin{hangparas}{1em}{1}

Brahe = Brahe, P. (1897). Per Brahe den äldres fortsättning af Peder Svarts krönika. Lund: Lunds Universitet.

Bur = Codex Bureanus. Ett fornsvenskt legendarium. (1847). Stephens, G. (ed.). Stockholm: Norstedt.

Gyll = Carl Carlsson Gyllenhielm: Egenhändige anteckningar rörande tiden 1597--1601. (1905). Almqvist, J. A. Sthlm: Historiska handlingar 20.

Hiarne = Hiärne, U. (1952). Stratonice. In M. von Platen (ed.). M von Platens utgåva Urban Hiärne: Stratonice. Stockholm: Norstedt.

Högström = Per Högström: Beskrifning öfwer Sweriges Lapmarker (1747). \url{https://project2.sol.lu.se/fornsvenska/} 

KM = Karl Magnus. 1887–1889. Prosadikter från Sveriges medeltid, ed. by Gustaf E. Klemming. Stockholm: Norstedt. Manuscript Holm D 4.

KS = Konungastyrelsen. (1984) Moberg, Lennart (ed.) (1984). Konungastyrelsen en filologisk undersökning. Uppsala: Svenska fornskriftsällskapet.

Linc = Linköpingslegendariet, Legenden om Sankt Amalberga. 1847. Ett fornsvenskt legendarium, ed. by George Stephens. Stockholm: Norstedt. Manuscript Linc B 70a.

Luc = Lucidarius. (1900). Geete, Robert (ed.) Svenska fornskriftsällskapets samlingar 33.

Pent = Pentateukparafrasen.1955. Fem Möseböcker på fornsvenska, ed. by Olof Thorell. Uppsala: Almqvist and Wiksell. Manuscript Codex Thott (Köpenhamn).

Spegel = Haqvin Spegel: Dagbok. (1923). Hildebrand, Sune (ed.). Stockholm: Norstedt.

ST = Själens tröst: Tio Guds bud förklarade genom legender, berättelser och exempel. 1873. Ed. by Gustaf E. Klemming. Stockholm: Norstedt. Manuscript Holm A108.

SVM = Sju vise mästare. (1887–89). In Prosadikter från Sveriges medeltid. Klemming, Gustaf E. (ed.). Stockholm: Norstedt.

Swart = Peder Swarts krönika (1912). Eden, Nils (ed.) Konung Gustaf I:s krönika, Stockholm.

Troj = Historia Trojana. (1892). In Historia Trojana. Geete, Robert (ed.). Stockholm: Norstedt.

All texts are available on Fornsvenska textbanken, at \url{https://project2.sol.lu.se/fornsvenska/}

\end{hangparas}




\sloppy
\printbibliography[heading=subbibliography,notkeyword=this]
\end{document}


